\chapter{How Peas Made Heredity Calculable}

\begin{kaichang}
Open your freezer: perhaps your family also has a bag of frozen peas. Put one in your palm: small, round, and green. Remember, it is not a ``food-factory component,'' but a ``child''---a seed produced after a pea plant was fertilized. Bury it in moist soil, and it can grow into a whole pea plant, flower, and produce another pod of peas.

Now make a prediction. Cross a plant producing round peas with one producing wrinkled peas---a plant ``marriage.'' What seeds will their children produce? Round? Wrinkled? ``Half-round, half-wrinkled,'' like the compromise color obtained by mixing two paints?

Write your guess on paper; we will return to it at the chapter's end. Over a hundred and sixty years ago, a monk named Mendel grew more than twenty thousand pea plants in a monastery garden, turning this question from ``reading appearances'' into ``arithmetic.'' Why did humble peas make heredity calculable for the first time? That is both our chapter title and the puzzle we will unpack.
\end{kaichang}

\section{Children Resemble Parents, but Are Not Blends}

First, give our subject an everyday name. A \textbf{trait} is a feature of an organism that you can ``see and count'': round or wrinkled pea seeds, yellow or green cotyledons, tall or short plants. In humans, examples include rolling your tongue into a tube, attached or free earlobes, bending the last thumb joint backward, and dimples when smiling. Check yourself in a mirror right now: no instruments are needed for these visible features.

Some traits are ``either-or'': you either can or cannot roll your tongue; few people ``half-roll'' it. Others are continuous: heights range from 1.4 meters to two meters, and skin has countless shades. This apparently minor distinction almost determined whether genetics could be invented. Remember it: Mendel was about to ``cheat'' by exploiting precisely this point.

In Mendel's time, the default theory of heredity sounded perfectly natural: \textbf{blending inheritance}. Each parent contributed half their ``qualities,'' like a drop of red ink mixed with a drop of blue, producing a purple child. Everyday experience seemed to fit: children's heights often lie between their parents', and their skin often seems an ``intermediate shade.''

But the theory hid a fatal prediction that few people had carefully followed through: if inheritance really blended, variation would be diluted by half each generation. Tall and short produce medium; medium and medium still produce medium. Within a few generations, everyone would cluster around the average, and variation would be ``stirred smooth'' until it vanished. Yet variation has never disappeared: average-height parents still have unusually tall or short children; two tongue-rollers can have a child who cannot roll; grandparents' features can skip a generation and suddenly ``revive'' in grandchildren.

So ``qualities'' are not ink: mixing does not erase them. But everyday observations alone will never persuade an opponent. ``More or less,'' ``it seems,'' and ``I think'' settle nothing. A decisive answer requires a countable experiment: choose an organism and several pairs of traits, mate the parents, and \textbf{count every plant}.

\section{Mendel's Garden: From Looking to Counting}

Gregor Mendel was born into a poor farming family in 1822 and later entered a monastery in Brünn, Austria (now Brno, Czechia). He was not a genius struck by a sudden flash: he formally studied mathematics and physics at the University of Vienna, bringing the physicist's skills of ``design controls, measure precisely, let numbers speak'' into the garden. Beginning in 1856, he spent about eight years doing something unprecedented. Every step was carefully designed and deserves a close look.

\textbf{Step one: why choose peas?} Peas seem almost made for this experiment. They self-pollinate: petals enclose a flower's stamens and pistil, so it naturally ``marries itself.'' This makes it easy to maintain purity across generations and develop stable ``pure lines.'' Yet their flowers are large enough to remove anthers by hand and apply another plant's pollen with a brush, giving the experimenter complete control of mating. They grow quickly, producing the next generation within a year. Each plant yields dozens or hundreds of seeds, providing plenty of offspring. Best of all, seeds themselves are the next generation: counting peas means counting ``children.''

\textbf{Step two: choose clear either-or traits.} Mendel did not try to study everything. He chose just seven sharply distinct pairs: round or wrinkled seeds, yellow or green cotyledons, purple or white flowers, inflated or constricted pods, green or yellow pods, axial or terminal flowers, and tall or short stems. He avoided continuous traits entirely. This was one of his smartest decisions: only clear either-or categories remove ambiguity from counting and allow ratios to emerge from the data.

\textbf{Step three: control mating.} He removed stamens before flowers matured (emasculation), covered flowers with paper bags to exclude a ``third party's'' pollen, and hand-pollinated them with pollen from the designated parent. Every offspring's parentage was recorded, with no loose entries in the ledger.

\textbf{Step four: count large numbers across generations.} Crossing pure-line parents produced the first offspring generation ($F_1$), whose seeds were all retained. Self-crossing $F_1$ produced the second generation ($F_2$), which he counted again. Over eight years, he examined more than twenty-eight thousand pea plants.

The first act was already surprising. Crossing pure round-seeded and wrinkled-seeded lines produced $F_1$ seeds that were \textbf{all round}. Wrinkling seemed swallowed whole---blending inheritance appeared to win a round, as one trait ``covered'' the other. But act two reversed the plot: self-crossing $F_1$ brought wrinkling back intact in $F_2$, in roughly one-quarter of the seeds. Mendel counted 5474 round and 1850 wrinkled, a ratio of about $2.96:1$. The other six pairs behaved alike: yellow to green cotyledons, $6022:2001$; purple to white flowers, $705:224$; tall to short stems, $787:277$. Every $F_2$ pair closely approached $3:1$.

Together, these observations unlock the answer. Wrinkling was not ``diluted'' in $F_1$, or it could not reappear intact in $F_2$ in a precise proportion. It was merely \textbf{masked} in $F_1$, preserved unchanged, then emerged a generation later. Heredity's carriers are not ink but discrete ``particles,'' passed along without being worn away.

\section{Particles: Seats for Invisible Factors}

Seeing a ratio is only the first step; explaining it is science. Mendel proposed a minimal internal picture of these ``particles,'' using just three assumptions:

\begin{itemize}[itemsep=2pt]
\item Each trait in a pea plant is governed by \textbf{a pair} of hereditary factors, one from the male parent and one from the female parent.
\item When gametes form (sperm cells in pollen and egg cells in ovules), the paired factors \textbf{separate}, with each gamete taking only one.
\item At fertilization, the two gametes each contribute one factor, forming a new pair \textbf{at random}.
\end{itemize}

These textbook statements seem unremarkable today, but they were revolutionary then. No one had dared propose discrete units inside organisms that came in pairs, separated, and recombined by a fresh draw of lots. Later scientists gave these units their modern names. The ``position'' determining a trait is a \textbf{gene}; different ``versions'' of a gene are \textbf{alleles}---round and wrinkled are two alleles of the same gene. An individual's combination of two alleles is its \textbf{genotype}; the appearance that combination ultimately produces is its \textbf{phenotype}. Two identical versions (round-round or wrinkled-wrinkled) are \textbf{homozygous}; different versions (one of each) are \textbf{heterozygous}.

Another concept needs immediate clarification: \textbf{dominance} and \textbf{recessiveness}. If only one version's effect appears in a heterozygote, the visible version is dominant (round), and the masked one recessive (wrinkled). Notice the wording: dominance and recessiveness describe ``what the heterozygote's phenotype looks like,'' nothing more. The factors do not fight until a winner takes charge; factors cannot fight. Dominance often reflects molecular interactions: perhaps one version makes a functional product while the other makes almost none, and one supply is enough, giving the heterozygote the dominant phenotype. A whole sequence of life processes connects genes to traits through chemicals such as proteins; Chapter 73 will discuss it. For now, remember: dominance concerns ``product supply and demand,'' not a ``ranking of factor strength.''

\section{Rules: Why $3:1$?}

With this particle-level picture, $3:1$ is no coincidence but an arithmetic problem. Let the round allele be $R$ and the wrinkled allele $r$. All $F_1$ individuals are heterozygotes, $Rr$. Half their gametes carry $R$ and half carry $r$, like tossing a fair coin. An $F_1$ self-cross is like tossing two coins together:

\begin{qiaoliang}
At each fertilization, the male gamete has probability $\frac{1}{2}$ of carrying $R$ and $\frac{1}{2}$ of carrying $r$; so does the female gamete. The two coin tosses are independent. As Chapter 8 explained, the probability of independent events occurring together is the product of their probabilities. Thus four combinations are equally likely:
\[
RR:\ \tfrac{1}{4},\qquad Rr:\ \tfrac{1}{4},\qquad rR:\ \tfrac{1}{4},\qquad rr:\ \tfrac{1}{4}.
\]
$Rr$ and $rR$ are the same genotype, giving the genotypic ratio $RR:Rr:rr = 1:2:1$. Since both $RR$ and $Rr$ have round phenotypes, the phenotypic ratio is
\[
\text{round}:\text{wrinkled} = \tfrac{3}{4}:\tfrac{1}{4} = 3:1.
\]
This is the \textbf{law of segregation}: paired alleles separate during gamete formation, entering different gametes. $3:1$ is not a peculiarity of peas, but the inevitable visible result of a ``$1:2:1$ genotype ratio'' plus ``dominant masking.''
\end{qiaoliang}

This logic immediately explains more numbers. Mendel next tracked two trait pairs together: round seeds with yellow cotyledons ($RRYY$) crossed with wrinkled seeds with green cotyledons ($rryy$) yielded all round-yellow $F_1$ ($RrYy$). Self-crossing produced four $F_2$ phenotypes---round-yellow, round-green, wrinkled-yellow, and wrinkled-green---in roughly a $9:3:3:1$ ratio. Its origin is clear: round to wrinkled is $3:1$, yellow to green is $3:1$, and the two pairs segregate independently, so their combined probabilities multiply. Expanding $(3+1)\times(3+1)$ gives exactly $9:3:3:1$. This is the \textbf{law of independent assortment} (free combination): different allele pairs segregate independently and combine freely. Remember its literal statement for now; Chapter 64 will show when it holds and when it fails. Genes are not always so ``unconnected.''

Why, you may ask, did the count give $2.96:1$, not exactly $3:1$? Because that is probability's nature. Four coin tosses commonly yield two heads and two tails, but three heads and one tail are hardly unusual. With four children, the probability of exactly two boys and two girls is only $\frac{3}{8}$. Smaller samples deviate more; larger samples are pressed toward the theoretical ratio by the law of large numbers. This is the deeper reason Mendel counted thousands of plants. Twenty-eight thousand over eight years averages about thirty-five hundred annually; spread over a growing season of a hundred-plus days, that meant pollinating, bagging, labeling, and observing dozens daily. The beautiful $3:1$ emerged from counting plants individually through eight whole springs.

Probability also works backward, answering questions that ``infer genotype from phenotype.'' Pick a round $F_2$ seed at random: what is the chance it conceals a wrinkled allele? Among round seeds, $RR:Rr$ is $1:2$, so the answer is $\frac{2}{3}$. This method of ``knowing the outcome and reasoning back to the cause'' is conditional probability. It may look like an arithmetic game, but it is fundamental to modern genetic counseling. We will use it again in the challenge section.

\section{Symbols: The Letters Finally Arrive}

Until now, we deliberately avoided all genetic symbols. Now they may appear, because you understand every entry in the ledger:

\begin{itemize}[itemsep=2pt]
\item A dominant allele takes a capital letter ($R$); its recessive counterpart takes the same lowercase letter ($r$). Case is purely a human bookkeeping convention, not an indication of size or strength.
\item Parents are $P$, the first offspring generation $F_1$, and the second $F_2$. ``$\times$'' means a cross; an ``$F_1$ self-cross'' means a plant pollinates itself.
\item Gamete types appear across an arrow: $Rr \rightarrow$ gametes $R$, $r$.
\end{itemize}

The handiest bookkeeping device is a \textbf{Punnett square}: list the two parents' gametes along the top and left, then enter each resulting zygote's genotype at their intersection. Every box has equal probability. Here is the square for $Rr \times Rr$:

\begin{table}[htbp]
\centering
\caption{A segregation Punnett square: each box represents $\frac{1}{4}$}
\begin{tabular}{c|cc}
\toprule
 & Gamete $R$ & Gamete $r$ \\
\midrule
Gamete $R$ & $RR$ (round) & $Rr$ (round) \\
Gamete $r$ & $Rr$ (round) & $rr$ (wrinkled) \\
\bottomrule
\end{tabular}
\end{table}

You can read the result at a glance: three round boxes, one wrinkled. Drawing the square matters much more than memorizing ``$3:1$.'' A memorized ratio may be forgotten; a square always lets you calculate it again.

A probability tree offers another bookkeeping method: ``$F_1$ produces gametes'' splits into two branches (each $\frac{1}{2}$), and each splits into two again. Multiply probabilities along the branches to obtain the zygote probabilities at the ends. Squares and trees are two sides of one coin: one uses area, the other paths, but both keep the same accounts. Two trait pairs require a $4\times4$ square with sixteen boxes (gametes $RY$, $Ry$, $rY$, $ry$). More boxes, unchanged logic. If you do not want to draw them, multiply: $(3:1)\times(3:1)=9:3:3:1$.

One final reminder: $R$ and $r$ are placeholders, not pictures of what alleles ``look like.'' Mendel knew neither their physical nature nor their location in cells, yet his arithmetic still worked. This is a rare scientific case of ``calculating the rule first and finding the object later.''

\section{Evidence: Answering the Skeptics}

\begin{zhengju}
A beautiful hypothesis is not enough for science: it must survive varied attempts to refute it. Mendel understood this and arranged a stringent test of his ``paired factors segregate and recombine'' hypothesis---the \textbf{testcross}.

A skeptic could offer an alternative: perhaps wrinkling was not ``hidden'' in $F_1$, but destroyed or transformed, with the wrinkled quarter in $F_2$ newly created. How could ``hidden'' be distinguished from ``destroyed''? Mendel crossed $F_1$ ($Rr$) with a pure wrinkled line ($rr$). If $F_1$ really retained an intact $r$, it should produce equal numbers of $R$ and $r$ gametes; combined with $rr$ gametes, the offspring should be half round and half wrinkled, or $1:1$. If $r$ had been destroyed or diluted, this ratio should not appear. The result closely matched $1:1$. Prediction confirmed: $F_1$ did conceal an unchanged wrinkling factor. The experiment's beauty is that it did not merely ``look'' at the same phenomenon again. It derived an \textbf{entirely new, previously unseen prediction} from the hypothesis, then let experiment vote.

What followed was not an inspiring success story. Mendel read his paper to the local natural history society in 1865, published it in 1866, and sent it to many prominent scholars---then heard almost nothing. Blending inheritance dominated; a monk applying mathematics to biology and writing about ``invisible hypothetical factors'' was either incomprehensible or unconvincing to most readers. The renowned botanist Nägeli reportedly suggested in his reply that Mendel ``study another plant'': hawkweed. Mendel did, with disastrous results. Hawkweed reproduces by apomixis, so its offspring did not follow segregation; years of work almost shook his confidence. Later, this became a regrettable historical episode: Mendel was not wrong; the recommended material happened not to fit his law.

The turning point came in 1900, sixteen years after Mendel's death. Three botanists---de Vries in the Netherlands, Correns in Germany, and Tschermak in Austria---independently performed crosses and reached similar conclusions, then found the paper that had slept for thirty-five years while searching the literature. All had to acknowledge Mendel's priority. Later that year, Bateson vigorously promoted the work in Britain, and the word ``genetics'' followed. A discovery could revive unchanged after thirty-five years of neglect not because of luck, but because its published numbers could be reproduced.

There is a still more intriguing postscript. In 1936, statistician Fisher reanalyzed Mendel's original data and found a ``happy problem'': the data were \textbf{too good}, fitting theoretical values more closely than natural statistical fluctuations permitted. Real experiments, probabilistically, should not look so tidy. Scholars still debate explanations: perhaps assistants intuitively excluded ``dubious-looking'' plants, or stopped counting once the expected ratio approached. But notice what is disputed: countless later crosses in peas and other organisms repeatedly reproduced $3:1$, leaving the law unshaken; the question concerns details of recordkeeping. This is how science works: it permits scrutiny of heroes' data because reproducibility, not heroic reputation, guarantees its laws.
\end{zhengju}

\section{Boundaries: When $3:1$ Fails}

Mendel's laws are powerful, but not the universe's universal grammar. Their conditions deserve to be remembered as explicit ``if … then …'' statements:

\begin{itemize}[itemsep=2pt]
\item \textbf{Diploidy, sexual reproduction, and nuclear genes} are needed for ``paired segregation.'' Bacteria have just one circular DNA molecule, so Mendelian ratios do not apply. Mitochondrial and chloroplast genes pass almost exclusively through the mother and also follow different rules (we will encounter them again around Chapter 66).
\item \textbf{Complete dominance} is needed for a $3:1$ $F_2$ phenotype ratio. If dominance is incomplete, the phenotype ratio changes---but note that the $1:2:1$ genotype ratio never changes; only its visible expression does (more next section).
\item \textbf{Random fertilization and comparable survival of genotypes}. If embryos of one genotype struggle to survive, the ratio immediately shifts. Some recessive alleles are lethal when homozygous, giving an ``odd'' $2:1$ $F_2$ ratio. This is not a counterexample, but a consequence of the law plus ``survival selection.''
\item \textbf{A sufficiently large sample}. Failure to see a recessive trait among four children tells you nothing; probability promises ratios only for large samples.
\item \textbf{Genuine independence between genes}. Independent assortment requires the gene pairs to lie on different chromosomes, or sufficiently far apart on one chromosome. Close neighbors pass along ``hand in hand,'' distorting $9:3:3:1$. This is linkage, discussed with chromosomes in Chapter 64. Interestingly, peas have exactly seven chromosome pairs, and most of Mendel's seven trait pairs lie on different chromosomes or far apart. His ``good luck'' concealed wise selection.
\item \textbf{A trait largely controlled by one gene pair}. Height, skin color, and crop yield involve dozens or hundreds of genes plus environmental effects; they do not yield $3:1$. These are Chapter 73's subjects.
\end{itemize}

Treat Mendel's laws as a precise map: within the territory of ``single genes, nuclear genes, large samples,'' they are exact. Beyond it, they do not become wildly wrong but change in explainable ways. Skilled users do not merely memorize the map: they know its boundaries.

\section{When Phenotypes Are Not Either-Or}

We said $3:1$ can ``change its face''; now look at those faces. These cases \textbf{extend}, rather than overturn, Mendel because the underlying genotypic segregation ratio remains unchanged.

\textbf{Incomplete dominance}. A cross between pure red- and white-flowered four-o'clocks (or snapdragons) produces neither red nor white $F_1$, but pink. Self-crossing $F_1$ yields red, pink, and white $F_2$ flowers in exactly $1:2:1$. What happened? The heterozygote makes just one supply of red-pigment material, so its color is ``halved.'' Neither allele masks the other; phenotype directly ``displays'' genotype. This is a gift to observers: for the first time, the phenotype ratio equals the genotype ratio, making Mendel's $1:2:1$ visible to the naked eye.

\textbf{Codominance}. In human ABO blood groups, the $I^A$ allele places A-type sugar chains on red-cell surfaces, while $I^B$ places B-type chains. An $I^A I^B$ heterozygote is not an ``A--B compromise'': both chain types appear \textbf{simultaneously}, giving type AB blood. Both versions are fully expressed, neither masking the other: codominance.

\textbf{Multiple alleles}. The same blood-group system has three circulating versions in the population: $I^A$, $I^B$, and $i$ ($i$ is recessive to both others). Watch the wording: multiple alleles describe diversity at the \textbf{population} level, not within one individual. Each person still has just two alleles at that locus, one from each parent. Type O is $ii$, AB is $I^A I^B$, and A may be $I^A I^A$ or $I^A i$. Six genotypes, four phenotypes: the same Punnett square can handle them all.

\begin{wuqu}
``Dominant alleles are stronger, more common, and more beneficial.'' This widespread claim is wrong in all three adjectives.

\textbf{More common?} Type O blood ($ii$) is recessive yet the commonest type in many world regions. Tongue rolling is dominant, yet non-rollers are plentiful. An allele's frequency depends on evolutionary history---migration, chance, selection---not whether it is dominant or recessive. Frequency is a population-level matter, which we will calculate more carefully around Chapter 85's discussion of evolution.

\textbf{More beneficial?} Huntington disease is a late-onset neurodegenerative condition whose disease allele is dominant: just one copy may lead to illness in middle age. Albinism and phenylketonuria, by contrast, are usually caused by recessive alleles. ``Dominant'' promises nothing about good or bad.

\textbf{Stronger?} Dominance describes the heterozygote's phenotype: molecular ``product supply and demand,'' not an arm-wrestling match between alleles. Moreover, an allele dominant in one tissue may behave differently in another environment or for another trait.

Remember: dominant $\neq$ powerful, dominant $\neq$ common, dominant $\neq$ superior. Next time an advertisement or rumor claims ``dominant genes are better,'' you have counterexamples ready.
\end{wuqu}

\begin{shiyan}
\textbf{Segregation in Your Palm} (about 30 minutes).

Materials: two one-yuan coins (or one red and one black card from each of two decks), paper, and a pen.

Steps: Name the coins ``male parent'' and ``female parent.'' Let heads represent dominant allele $R$ and tails recessive $r$, making each coin a heterozygous parent $Rr$'s ``gamete generator.'' Toss both together and record: two heads as $RR$, one head and one tail as $Rr$ (order does not matter), two tails as $rr$. Count $RR$ and $Rr$ as ``dominant phenotype'' and $rr$ as ``recessive phenotype.'' Repeat 40 times and calculate genotype and phenotype ratios. Then pool data with your deskmate to reach 80 and 160 trials, recalculating each time.

Observations: After only 12 tosses, your result may be far from $1:2:1$. More tosses press the ratio toward theory. Your technique is not becoming more accurate: the law of large numbers is unfolding in your palm. Mendel's twenty-eight thousand plants tackled precisely the same luck.

Safety: Keep small objects such as coins and beans away from young children and pets to prevent swallowing. If you use beans, do not eat those repeatedly handled or dropped on the floor. One serious final reminder: this game models an abstract probability mechanism. Do not use it to ``tell fortunes'' about relatives' traits or judge anyone.
\end{shiyan}

\section{Back to the Frozen Peas}

Return to our opening pea. You can now tell yourself its whole story.

Why do pure round and pure wrinkled parents not produce ``half-round, half-wrinkled'' offspring? Because traits are particles, not paint. $F_1$ has genotype $Rr$; the round allele supplies enough product, masking the wrinkled version, so the entire pod of $F_1$ seeds is round. Why does wrinkling ``revive'' in $F_2$? It never disappeared: preserved intact in $F_1$, it segregates and recombines until one-quarter of zygotes receive $rr$, and wrinkling returns openly. If your written guess agrees, congratulations on independently reaching 1865. If not, better still: now you know why.

A bonus quietly reconnects this chapter to the book's chemical thread. Why are wrinkled peas wrinkled? Modern molecular biology found that an extra sequence inserted into the wrinkled allele disables an enzyme that catalyzes starch branching. The seeds have more sugar and less starch, then wrinkle as they lose water during ripening. Follow dominance to its roots, and you reach an enzyme's working state. Differences in hereditary information ultimately become \textbf{differences in chemical composition}. This hidden thread of ``information becoming matter'' unfolds fully in Chapters 66--69.

Incidentally, commercially sold sweet, crisp peas often have wrinkled ancestry: their sweetness comes from more sugar and less starch. The sweetness in your fried rice is the legacy of that quarter of Mendel's $F_2$.

\section{Further: Breeding Fields and Hospitals}

Mendel's arithmetic left the garden long ago.

In breeding fields, it predicts outcomes. When crossing two parents with different strengths, breeders know that dominant traits appear in $F_1$, while masked recessive advantages emerge in roughly one-quarter of $F_2$. They therefore need a large population to find the desired combination. The ledgers behind hybrid rice and disease-resistant wheat bear the shadows of $3:1$ and $9:3:3:1$.

In hospitals, it calculates risk. For conditions usually caused by recessive alleles, such as phenylketonuria and albinism, carriers ($Aa$) are themselves completely healthy. If both parents happen to be carriers, each child has about a $\frac{1}{4}$ chance of being affected and a $\frac{1}{2}$ chance of being a carrier like the parents. Genetic counselors apply a serious version of this chapter's conditional probability. But respect the book's usual boundary: these are \textbf{population-level} probabilities. For a particular family, inheritance gives no verdict of ``will'' or ``will not,'' only a reference for ``how likely.'' Consult qualified professionals at a medical institution if you have real concerns; this book provides no individual diagnosis or treatment advice.

Finally, take this chapter's greatest mystery with you. Mendel calculated his factors precisely but died without knowing \textbf{where they lived or what they were made of}. His $R$ and $r$ were ledger symbols without physical objects. Around the time of his death, rapidly improving microscopes let cytologists see threads dancing inside dividing nuclei: chromosomes. They occur in pairs, separate in division, halve during gamete formation, and restore their pairs at fertilization---a dance strikingly parallel to Mendel's segregation and recombination. Coincidence, or do factors ride chromosomes? If so, how can two genes on the same chromosome ``assort independently''? Next chapter, we zoom into the nucleus to see chromosomes bring Mendel's arithmetic into cells---and discover where it stumbles.

\begin{tiaozhan}
\textbf{Reasoning Backward: From Phenotype to Genotype.} A rare disease caused by recessive allele $a$ affects only $aa$ individuals. A healthy couple each have an affected sibling, with no information available about either side's grandparents. What is the probability that their first child will be affected?

Break it down. An affected sibling ($aa$) means that person's parents must both be carriers, $Aa \times Aa$. The person is healthy, excluding $aa$. Conditional on being ``healthy,'' the probabilities are $\frac{1}{3}$ for $AA$ and $\frac{2}{3}$ for $Aa$ (the same structure as the round seed's $\frac{2}{3}$). The partners are independent events, so the probability both are carriers is $\frac{2}{3}\times\frac{2}{3}=\frac{4}{9}$. Two carriers have probability $\frac{1}{4}$ of an affected child. Multiplying gives
\[
\frac{2}{3}\times\frac{2}{3}\times\frac{1}{4}=\frac{1}{9}\approx 11\%.
\]
Every step is just conditional probability. The difficulty is not arithmetic but deciding ``which option to exclude at which step.'' This is the core algorithm of genetic counseling. You may skip this section without losing the main thread.
\end{tiaozhan}

\textbf{Translator safety note:} The challenge's numerical answer assumes both parents of each partner are unaffected carriers; an affected sibling alone does not establish that assumption. Do not use this simplified exercise to assess a real family's risk. A genetics professional considers the diagnosis, inheritance pattern, and family information; see \href{https://medlineplus.gov/genetics/understanding/inheritance/riskassessment/}{MedlinePlus genetic-risk guidance}.

\section*{Try It Yourself}

\begin{sikaoti}
\item Testcross a heterozygous tall pea ($Tt$) with a short pea ($tt$). Draw the Punnett square, predict genotype and phenotype ratios, and explain in one sentence why the experiment can ``see'' what is hidden inside $F_1$.
\item An $Rr$ self-cross yields only 4 seeds. What is the probability of exactly 3 round and 1 wrinkled? (Hint: draw a probability tree or list the four positions; each seed has probability $\frac{3}{4}$ of being round.) How far is the answer from $100\%$? Why does this show Mendel could not grow only four plants?
\item Crossing red ($RR$) and white ($rr$) four-o'clocks gives pink $F_1$, which self-cross. Draw the square and give the $F_2$ genotype and phenotype ratios. Why does the phenotype ratio here ``faithfully display'' the genotype ratio for the first time?
\item Draw an information-flow diagram from ``paired alleles in parental body cells'' through ``gametes'' to ``zygote'' and ``phenotype.'' Mark where segregation occurs, where random combination occurs, and where dominance or recessiveness becomes visible.
\item A type A father and type B mother have a type O child. Write the parents' possible genotypes and use a square to show where the O child comes from. What is the probability that their next child is type AB?
\item Your deskmate says: ``Tongue rolling is dominant, so after a few dozen generations everyone will roll their tongue.'' Identify the two concepts confused in this reasoning, and give two counterexamples from the chapter.
\end{sikaoti}
